\phantomsection
\addcontentsline{toc}{chapter}{Tóm tắt đồ án}
\thispagestyle{empty}
% Doi ma mau HEX tai day (vi du: C62828, 1E88E5, 2E7D32)
\definecolor{TomTatTitleColor}{HTML}{880000}

\begin{center}
    {\Large\textbf{\textit{\textcolor{TomTatTitleColor}{TÓM TẮT ĐỒ ÁN}}}}
\end{center}

\vspace{1.5em}

Đồ án tập trung nghiên cứu và thiết kế mô đun tải trọng thí nghiệm đa năng cho vệ tinh nhỏ. Mục tiêu chính của đề tài là xây dựng một kiến trúc phần cứng và phần mềm có khả năng điều khiển thí nghiệm, giao tiếp với các board chức năng, thu nhận dữ liệu và hỗ trợ mở rộng cho nhiều cấu hình tải trọng khác nhau.

Hệ thống được tổ chức theo kiến trúc mô đun, trong đó board EXP đóng vai trò điều khiển trung tâm, quản lý trình tự thí nghiệm, xử lý dữ liệu và phối hợp hoạt động giữa các khối ngoại vi. Board IF đóng vai trò giao tiếp mở rộng, hỗ trợ chuẩn hóa kết nối giữa EXP với các thành phần khác trong hệ thống. Cách tổ chức này giúp tách biệt phần điều khiển lõi với phần giao tiếp ngoại vi, từ đó tăng tính linh hoạt khi thay đổi hoặc mở rộng cấu hình thí nghiệm.

Trong phạm vi đồ án, cụm board LASER và PHOTO được sử dụng như một cấu hình thử nghiệm nhằm kiểm chứng khả năng hoạt động của mô đun. Board LASER đảm nhiệm chức năng phát tín hiệu quang, trong khi board PHOTO thu nhận tín hiệu sau khi tương tác với mẫu thí nghiệm và chuyển đổi thành dữ liệu phục vụ xử lý. Thông qua cấu hình này, đồ án đánh giá khả năng điều khiển ngoại vi, thu nhận tín hiệu và truyền dữ liệu của hệ thống.

Kết quả thực hiện của đồ án bao gồm thiết kế phần cứng các board chức năng, xây dựng phần mềm nhúng phục vụ điều khiển và truyền nhận dữ liệu, đồng thời kiểm thử hoạt động của hệ thống theo các kịch bản thí nghiệm. Đề tài hướng đến việc xây dựng một nền tảng tải trọng thí nghiệm có tính mô đun, có khả năng kế thừa và phát triển cho các ứng dụng nghiên cứu trên vệ tinh nhỏ trong tương lai.

\vfill
